\begin{lemma}[Compact-$p$ interface]\label{lem:compact-p}
For every $n\geq2$ and finite $P\geq3$, there are common
$\theta(n,P)>0$ and $C(n,P)<\infty$ such that every bounded weak
$p$-harmonic solution on $B_1$, $3\leq p\leq P$, satisfies
\[
 \|Du\|_{L^\infty(B_{1/2})}+[Du]_{C^{0,\theta}(B_{1/2})}
 \leq C(n,P)\osc_{B_1}u.
\]
This includes zero gradients. An entire weak solution for one fixed
$p\geq3$ with bounded gradient is affine.
\end{lemma}
\begin{proof}
We give the common choices in \cite[Sections 1--5]{Uhlenbeck}; a
fixed-$p$ estimate alone would not establish uniformity. Work first
in dimension $m\geq3$, and distinguish Uhlenbeck's growth exponent
$s=(p-2)/2\in[1/2,S]$, $S=(P-2)/2$, from the present $p$. Set
$\omega=Du$, $Q=|\omega|^2$, $\rho(Q)=Q^s$, and set her structural
constant $C$ to zero. The gradient/curl complex satisfies her (1.2),
and $\omega$ solves her system (1.1).

\paragraph{Common structural and preliminary estimates.}
Fix
\[
 a=\min\{1/2,1/(m-2)\},\qquad
 K_*=2+(1+2S)+S\max\{1,S\}.
\]
The exponent $a$ is positive, at most every $s$, and strictly below
$2/(m-2)$, so $2(1+a)$ is one fixed admissible Sobolev exponent.
The ellipticity expression is
$\rho+2Q\rho'=(1+2s)Q^s$, and
\begin{equation}\label{eq:compact-structural}
 |s x^s-s y^s|\leq S\max\{1,S\}(x+y)^{s-a}|x-y|^a
 \qquad(x,y\geq0).
\end{equation}
Indeed divide by $(x+y)^s$ when $x+y>0$. On $[0,1]$, the power
$t^s$ is $s$-H\"older with constant one for $s\leq1$, and is
$S$-Lipschitz for $s\geq1$. Since $a\leq s$ and $a\leq1$, either
bound implies \eqref{eq:compact-structural}; the zero case is direct.
Thus (1.3), (1.3)$'$ and (1.4) of the source hold with $K_*$.
Differentiability of $\rho$ is required only at $Q>0$.

The source's (1.7), (1.10) and Sections 2--3 apply to
\[
 H_s(Q)=\frac{s+1/2}{s+1}Q^{s+1}
       =\frac{p-1}{p}|Du|^p.
\]
Its prefactor is in $[2/3,1]$. The divergence-form operator in
(1.10) has common measurable ellipticity bounds controlled by
$K_*$. Fix a mean-value exponent $\beta$ strictly between $1$ and
$m/(m-2)$ in (2.1), and fix the Sobolev exponent. Moreover,
\[
 |D\sqrt{H_s(Q)}|\leq(S+1)Q^{s/2}|D\omega|.
\]
The weighted differentiability and energy estimate (3.1), followed
by these fixed Sobolev and Moser estimates in the proof of (3.2),
give one $C_2(m,P)$ such that
\begin{equation}\label{eq:compact-moser}
 \sup_{B_r}H_s(Q)\leq C_2(m,P)r^{-m}
                         \int_{B_{4r}}Q^{s+1}.
\end{equation}
The constants here and in (3.3) involve $K_*$, $(S+1)^2$, fixed
cutoff factors, Sobolev and Poincar\'e constants, and Moser's
$c(\beta)$ and weak-Harnack constant $c'(1)$. None requires a
parameter-dependent threshold choice. In the scalar case the
potential for $\omega$ is $u$ itself; no varying-exponent Hodge
potential estimate is implicit. The weak differentiability and
subsolution justification in Sections 2--3 includes $Q=0$.

\paragraph{One common small-excess iteration.}
Fix the relative reference size $\eta=1/2$ in the source's (4.1).
Normalize $\sup Q=M=1$ by homogeneity; $M=0$ is trivial. During the
iteration maintain $|w_i|\leq1$ and $|w_i|^2\geq1/4$. The
normalized constant-coefficient linearization at a reference $w$ is
\[
 I+2s\frac{w}{|w|}\otimes\frac{w}{|w|},
\]
with eigenvalues in $[1,P-1]$. Its comparison energy and interior
estimates in (4.2) therefore have common constants.

Here is a uniform derivation of the nonlinear error estimate,
without leaving the power-map constant of printed (4.6) implicit.
Put
\[
 A_s(v)=|v|^{2s}v,\qquad F_s(v)=|v|^s v,\qquad
 G_s(v,w)=\frac{A_s(v)-A_s(w)-DA_s(w)(v-w)}{\rho(|w|^2)}.
\]
For $q=2s\in[1,2S]$, the second derivative of $A_s$ away from zero
is bounded by $C_mq(q+3)|v|^{q-1}$. Integration on a segment,
split at zero when necessary, shows that $DA_s$ is uniformly
Lipschitz on the unit ball. Since $\rho(|w|^2)\geq4^{-S}$,
\begin{equation}\label{eq:compact-remainder}
 |G_s(v,w)|\leq C(m,P)|v-w|^{1+a}
 \quad (|v|,|w|\leq1,\ |w|^2\geq1/4).
\end{equation}
To invert the power map uniformly, let $R=\max\{|v|,|w|\}$.
The inequality $DF_s(z)\geq|z|^sI$, integrated over the quarter of
the joining segment nearest a maximal endpoint, gives
\[
 (F_s(v)-F_s(w))\cdot(v-w)
       \geq2^{-S-2}R^s|v-w|^2.
\]
Here that segment has norm at least $R/2$, while $R\geq1/2$.
Cauchy--Schwarz consequently yields
\begin{equation}\label{eq:compact-inverse}
 |v-w|\leq2^{2S+2}|F_s(v)-F_s(w)|.
\end{equation}

Let $E=\int|\omega-w|^2$ on a fixed outer ball. The weighted
estimate (4.8) has a common scalar constant: $DA_s(\omega)$ has
eigenvalues between $Q^s$ and $(P-1)Q^s$, so testing the
differentiated equations against cutoff squared times
$\omega_j-w_j$ gives a weighted-energy factor at most
$4(P-1)^2$ times fixed cutoff constants. The source's
difference-quotient argument justifies this at weak solutions and
zeros of $\omega$. Also
\[
 \|DF_s(v)\|\leq(S+1)|v|^s,\qquad
 |F_s(v)-F_s(w)|\leq(S+1)|v-w|
 \quad (|v|,|w|\leq1).
\]
Thus on nested fixed balls,
\[
 \|F_s(\omega)-F_s(w)\|_{W^{1,2}}^2\leq C(m,P)E.
\]
The fixed embedding into $L^{2(1+a)}$, together with
\eqref{eq:compact-remainder}--\eqref{eq:compact-inverse} and the
linear comparison tested against its zero-boundary difference,
shows that the squared $L^2$ error of the linearized comparison
field is at most $C(m,P)E^{1+a}$, as required in (4.3)--(4.4).

Specify the next reference as follows. If $\widetilde\omega_i$ is
the linearized comparison to $\omega^i-w_i$ on the fixed inner
ball, take
\[
 \widehat w_{i+1}=w_i+\widetilde\omega_i(0),\qquad
 w_{i+1}=\operatorname{proj}_{\overline B_1}(\widehat w_{i+1}).
\]
The center value is controlled by the common interior estimate;
its gradient bound gives the $O(t^2)$ excess improvement. Projection
does not increase the error to any $\omega^i(x)\in\overline B_1$
or the distance to $w_i\in\overline B_1$. With the fixed half-ball
rescalings in (4.9)--(4.11) absorbed into one $C_3=C_3(m,P)\geq1$,
let $E_i$ be the $L^2$ excess on the rescaled unit ball and
$\omega^i$ the corresponding rescaled field. We obtain
\begin{equation}\label{eq:compact-iteration}
 E_{i+1}\leq C_3(t^2+t^{-m}E_i^a)E_i,\qquad
 |w_{i+1}-w_i|\leq C_3\sqrt{E_i}.
\end{equation}
Choose once
\begin{equation}\label{eq:compact-choices}
 \begin{split}
 t&=\min\{1/8,1/(4C_3)\},\\
 \epsilon&=\min\left\{1,\frac1{256C_3^2},
              \left(\frac{t^{m+1}}{4C_3}\right)^{1/a}\right\}.
 \end{split}
\end{equation}
If $E_0\leq\epsilon$, then $E_{i+1}\leq(t/2)E_i\leq tE_i$,
and
\[
 \sum_i|w_{i+1}-w_i|
 \leq\frac{C_3\sqrt\epsilon}{1-\sqrt t}
 \leq2C_3\sqrt\epsilon\leq1/8.
\]
All references have norm at most one, so their squared norms differ
from $|w_0|^2\geq1/2$ by at most twice this drift. Hence
$|w_i|^2\geq1/4$, closing the induction. The source's
(4.12)--(4.15), with their common geometric-series ratio, now give
(4.1) with one threshold $\epsilon_*>0$ and one $1/2$-H\"older
constant $K_4(m,P)$. Fixed unit-ball versus $B_2/B_1$ conventions
only decrease $\epsilon_*$ by fixed radius and volume factors.

\paragraph{One common decrease factor, including critical points.}
Proposition (5.1) uses (3.3), the fixed weak-Harnack constant
$c'(1)$, Poincar\'e for $F_s(\omega)$, and the factors
$(s+1)^2$ and $2^{2s+2}$. These are bounded by $(S+1)^2$ and
$2^{2S+2}$; all other factors are the common $K_*$, cutoff,
Sobolev and Moser constants just listed. The radial maps selecting
the reference have exponent $s+1\in[3/2,S+1]$, so neither their
exponents nor their reciprocals introduce a degenerating choice.
Choose one $C_5(m,P)\geq1$ for that proposition. The two required
inequalities on printed p.~238 are
\[
 1-C_5\sqrt\lambda\geq2^{-(s+1)/2},\qquad
 \lambda C_5 2^{s+1}\leq\epsilon_*.
\]
They hold simultaneously for all $s$ with
\[
 \lambda_*=\min\left\{1/4,\frac1{16C_5^2},
                       \frac{\epsilon_*}{2C_5 2^{S+1}}\right\}>0.
\]
Indeed the first left side is at least $3/4$, while its right side
is at most $2^{-3/4}<3/4$; the second inequality is direct. Set
\[
 \theta=-\frac{\log(1-\lambda_*)}{2\log4}>0,
 \qquad\sqrt{1-\lambda_*}=4^{-\theta}.
\]
Then $\theta<1/2$. Repeated maximum decrease in (5.4)--(5.5)
gives the $r^\theta$ modulus; at the first scale where the
perturbation alternative applies, the common $1/2$-H\"older bound
is stronger. If the decrease continues forever, the maximum tends
to zero, covering a critical gradient without linearizing at zero.
Consequently
\begin{equation}\label{eq:compact-local-holder}
 [\omega]_{C^{0,\theta}(B_{r/2})}
       \leq K_0(m,P)r^{-\theta}\sqrt M
 \quad\text{if }\sup_{B_r}|\omega|^2\leq M.
\end{equation}
Translation supplies the estimate at every center.

\paragraph{The oscillation bound and dimension two.}
Normalize $0\leq u\leq1$ by subtracting its essential infimum and
dividing by its essential oscillation; zero oscillation is trivial.
A cutoff $\chi=1$ on $B_{7/8}$, supported in $B_1$, with
$|D\chi|\leq16$, is used in the admissible weak test $\chi^pu$.
H\"older's inequality, with the zero-energy case separated, gives
\[
 \int\chi^p|Du|^p\leq
 p\left(\int\chi^p|Du|^p\right)^{(p-1)/p}
   \left(\int|u|^p|D\chi|^p\right)^{1/p}.
\]
Sobolev approximation justifies the test. Therefore
\[
 \int_{B_{7/8}}|Du|^p\leq(16P)^P|B_1|=:E_*.
\]
Apply \eqref{eq:compact-moser} with centers in $B_{3/4}$ and
$r=1/32$, so $B_{4r}$ stays in $B_{7/8}$. It gives
\[
 \sup_{B_{3/4}}|Du|^p\leq(3/2)C_2 32^m E_*=:A_*,\qquad
 \|Du\|_{L^\infty(B_{3/4})}\leq\max\{1,A_*\}^{1/3}=:L_*.
\]
Use \eqref{eq:compact-local-holder} on $B(x,1/4)$ for
$x\in B_{1/2}$. Pairs separated by at most $1/8$ obey that local
bound; more distant pairs obey
$|Du(x)-Du(y)|\leq2\,8^\theta L_*|x-y|^\theta$. This proves the
claimed gradient norm and full seminorm after restoring oscillation.
For original dimension two, extend $u$ constantly in a third
coordinate. Fubini shows it is an actual weak $p$-harmonic function
on $B_1^2\times\R$, with the same gradient norm and oscillation.
Apply the three-dimensional estimate on $B_1^3$ and restrict back.
This constant extension is distinct from the affine lift $u+b$.
Finally, for a single fixed $p$, scale
\eqref{eq:compact-local-holder} on arbitrarily large balls for an
entire bounded-gradient solution. Its gradient oscillation on any
fixed ball tends to zero, so it is affine, also as in the source's
Bernstein consequence (6.1).
\end{proof}
